% Response to the Google Deep Review — Round 2
% For "Trustworthy Agentic Automation for Intent-Based Self-Healing in
% Software-Defined Networks: Three Deterministic Gates around a Bounded Agent"
%
% Compile with:  tectonic REVISION_REPORT_ROUND2.tex
\documentclass[11pt,a4paper]{article}
\usepackage[margin=1.9cm]{geometry}
\usepackage{mathptmx}
\usepackage{amsmath,amssymb}
\usepackage{booktabs,longtable,array,colortbl}
\usepackage{enumitem}
\usepackage{xcolor}
\usepackage{fancyhdr,titlesec}
\usepackage[hidelinks]{hyperref}
\usepackage{microtype}
\emergencystretch=2.5em
\sloppy

\definecolor{ink}{HTML}{1F2A44}
\definecolor{accent}{HTML}{0F5EA8}
\definecolor{quote}{HTML}{5A6B7C}

\hypersetup{
  pdftitle={Response to the Google Deep Review — Round 2},
  pdfauthor={The author},
  colorlinks=true, urlcolor=accent, linkcolor=accent, citecolor=accent}

\titleformat{\section}{\Large\bfseries\color{ink}}{\thesection}{1em}{}
\titleformat{\subsection}{\large\bfseries\color{ink}}{\thesubsection}{0.8em}{}
\setlist{nosep,topsep=3pt}

\pagestyle{fancy}
\fancyhead[L]{\small\color{quote}Response to the Google Deep Review --- Round 2}
\fancyhead[R]{\small\color{quote}\thepage}
\fancyfoot{}

\newenvironment{reviewer}{\par\medskip\noindent\begin{minipage}{\linewidth}\small\itshape\color{quote}\textbf{Reviewer:}\par}{\end{minipage}\medskip\par}
\newenvironment{response}{\par\noindent\begin{minipage}{\linewidth}\small\textbf{Response.}\par}{\end{minipage}\medskip\par}

\begin{document}

\begin{titlepage}
\vspace*{2.8cm}
{\color{ink}\Huge\bfseries Response to the Google Deep Review\\[0.3em]\bigskip}
{\color{ink}\Large Round 2 Report\bigskip}
{\color{ink}\large\bfseries
\emph{Trustworthy Agentic Automation for Intent-Based Self-Healing in
Software-Defined Networks: Three Deterministic Gates around a Bounded Agent}\\[0.6em]}
{\large Revised submission, November 2026}

\vspace{1.2cm}
{\small
\begin{tabular}{@{}ll@{}}
\toprule
\textbf{Round 2 review} & \texttt{final\_submission/GA\_DEEP\_REVIEW\_ROUND2.pdf} \\
\textbf{Review source} & \texttt{GA\_DEEP\_REVIEW\_ROUND2.md} (byte-identical to the pasted text) \\
\textbf{Revised paper} & \texttt{final\_submission/paper/paper.pdf} \\
\textbf{Round 1 letter} & \texttt{final\_submission/REVISION\_REPORT.pdf} \\
\textbf{One command} & \texttt{bash final\_submission/scripts/reproduce.sh} \\
\bottomrule
\end{tabular}
}

\vspace{1.5cm}
\noindent\textbf{A note on how to read this document.}

Every Round 2 comment is reproduced (in the reviewer's words, lightly
trimmed, never paraphrased beyond recognition) and is followed by a
numbered response stating exactly what changed and where to check it. The
review's own structure is kept: Sections~1–5 mirror the review's
\emph{Potential Mistakes and Improvements} and \emph{Minor Corrections and
Typos} groupings (introduction/related work and framing; system model and
three gates; experimental setup and ablations; limitations and roadmap), and
Section~6 covers the citation verification. Where a claim was removed
rather than fixed, the removal is stated plainly. Where the reviewer was
right and the artifact was wrong, the defect is described and the fixing
test or measurement is named. Where the review overstated, we say so and
why.

\vfill
\end{titlepage}

\section{Summary of changes, and what the review actually found}

The Round 2 review endorsed the artifact's structure (the three gates, the
terminal-outcome taxonomy, the truth audit) and isolated four genuine
problems in the rigor of the evaluation and the framing:

\begin{enumerate}
  \item The paper hedged ``agentic LLM automation'' while evaluating a
    deterministic keyword-parser fallback with a deterministic planner; the
    gating architecture (not the model) was what the evidence supported.
  \item The offered-load scenario (S7) sat below the utilisation threshold
    it was meant to exercise, and conflated offered load with the delay
    fault it shares with S2.
  \item The active rerouting scenarios S3, S7 and S8 were single trials,
    while S2 was reported as a distribution.
  \item Several internal contradictions remained: the interaction of the
    loss clause in Eq.~(4) with the verdict of Eq.~(7), the placement of
    the improvement test (\texttt{diagnose} vs \texttt{validate}), the
    ordering of freshness vs.\ verdict checks, and roadmap items that
    re-announced already-completed fixes as future work.
\end{enumerate}

\noindent All four were fixed. In addition, the review's recommended
additions to related work and to the roadmap were made (Section~2 and
Section~5), the citation verification's two flagged entries were corrected
(Section~6), and every minor correction was applied. The paper's measured
results were regenerated end-to-end: S2, S3, S7, S8 now each report five
independent trials, and the wall-clock duration of the closed loop is
reported for the active-recovery runs (Section~4). No number in the paper is
hand-typed: every tabular value is generated from the recorded data by
\texttt{make\_paper\_tables.py}, and the paper rebuilds with \texttt{tectonic}.

\subsection{Ground rule: what this letter will not claim}

The evaluation is unchanged in kind: no model endpoint is exercised in any
evaluated run. Every claim about the gating architecture is backed by a
recorded run or a test; every claim about the untested model boundary is
explicitly labelled future work. The reviewer's \S4.4 concerns (Gate 2's
rejection path and the $M \rightarrow D$ boundary are unexercised in the
data) are acknowledged as scope limits rather than papered over---we state
them plainly in Sections~VI-C, X and XI rather than inventing scenarios to
fill them.

\section{Framing and related work}

\subsection{Alignment between the ``Agentic / LLM'' framing and the evaluated
deterministic system (Title, Abstract, Sections I, II-D \& XI)}

\begin{reviewer}
Candidate reroutes are proposed strictly by a deterministic graph planner\dots
Gates 1, 2, and 3 do not actually evaluate or filter LLM action proposals.
\dots Evaluating without an active LLM leaves the $M \to D$ trust boundary
\dots untested. \dots readers may interpret this concluding lesson as an
empirically validated finding about LLM behavior.
\end{reviewer}

\begin{response}
Three changes were made, exactly matching the suggested improvements
(a)--(c):

(a) The Abstract now states that the candidate reroute is proposed by a
deterministic graph planner computed from the live topology, that ``the
evaluation tests that gating architecture in isolation, with no model
configured and its deterministic keyword-parser fallback in use,'' and that
the behaviour of the $M \rightarrow D$ boundary against a live, noisy or
adversarial model is future work, not a result of this paper. The Abstract
also no longer claims ``plateaus below line rate'' as if it were a
fundamental property; it states the offered-load scenario sustains a
narrow 60 Mbit/s bottleneck across the 80\% threshold on live byte
counters.

(b) The Abstract and Section~I now say that the three gates stand between
the \emph{planner's candidate action} and the data plane, not between
``the model's suggestion''. Section~I adds the sentence: ``the gates do
not judge model outputs; they judge candidate actions.'' The Section~IV-C
trust-boundary framing was already correct and is unchanged.

(c) The Conclusion's ``design lesson'' paragraph is now explicitly framed
as an architectural design principle---what was measured is the gating
architecture in isolation, with a deterministic parser and a deterministic
planner---and it says in so many words that we ``did not measure how the
boundary behaves when the untrusted component is a live, noisy or
adversarial model.'' Roadmap item~1 (``Exercise the active model'')
remains the obvious next experiment and now cites NetConfEval
(Changjie Wang et al., NetConfEval, Proc.~ACM Netw., 2024) and NIKA (Zhihao Wang et al., NIKA, 2025) as the concrete
grounding. No claim about LLM behaviour appears anywhere in the revised
paper.
\end{response}

\subsection{Positioning relative to classical closed-loop SDN intent
reconciliation systems (Sections II-B \& II-C)}

\begin{reviewer}
Contrasting dynamic post-action latency verification primarily against
static header-space/invariant checkers \ldots leaves a gap regarding prior
closed-loop SDN intent monitoring and rerouting frameworks \ldots as well as
structured natural-language IBN translation pipelines.
\end{reviewer}

\begin{response}
Section~II-C gained a paragraph positioning the gates against the closest
prior art, all verified against the publishers' own records: ONOS's Intent
Monitor and Reroute service (Davide Sanvito et al., ONOS Intent Monitor and Reroute, NetSoft, 2018) (monitor and reroute;
no post-change measured outcome verdict and no refusal semantics); Lumi
(Arthur S. Jacobs et al., Hey, Lumi!, USENIX ATC, 2021) (natural-language intents into ONOS configurations;
checks at configuration time); NetConfEval (Changjie Wang et al., NetConfEval, Proc.~ACM Netw., 2024)
(benchmarks for LLM configuration); Confucius (Zhaodong Wang et al., Confucius, SIGCOMM, 2025) (a
production multi-agent LLM framework at Meta, models propose and
domain-specific tools check); and Protogeros et al.\ (Ioannis Protogeros et al., Benchmarking LLM-Driven Network Configuration Repair, arXiv:2604.22513, 2026)
(verifier-in-the-loop repair benchmarks). Section~II-D now also cites
NetConfEval, Confucius and Protogeros as verifier-gated frameworks. The
claim structure of the whole positioning is explicit: prior work halts at
configuration-time correctness or acts inside a model-centred loop; this
paper's gates judge the \emph{measured runtime outcome} and make refusal a
first-class, audited outcome.
\end{response}

\subsection{Minor correction: Figure 1 caption terminology}

\begin{reviewer}
The caption states ``The cyan arrow is the protected video intent
$h1\to h4$.'' \ldots the video flow is the controlled (target) intent
flow, whereas the voice flow is the protected flow.
\end{reviewer}

\begin{response}
Fixed: the caption now reads ``The cyan arrow is the controlled video
intent $h1\!\to\!h4$; the protected voice flow that a recovery must not
harm is $h3\!\to\!h4$.''
\end{response}

\section{System model, mathematical formulation and the three gates}

\subsection{Inconsistent responsibility and ordering between candidate
pruning (\texttt{diagnose}) and safety validation (\texttt{validate})}

\begin{reviewer}
None of the nine safety validation checks \dots is an ``improvement test''.
\dots If candidates that fail to improve upon the current path are already
dropped during diagnose (Line 19) before reaching validate (Line 20), then
\dots $\mathcal{P}_c$ is already empty \dots
\end{reviewer}

\begin{response}
The paper now describes the code as it is, and Algorithm~1 was rewritten to
the implementation's real order instead of to the old pseudocode's:

\texttt{diagnose} ranks suspect links (see Section~3.3) and enumerates
candidate paths, dropping only the candidate identical to the path the flow
is already on (installed as the Eq. (6) ranking's own semantics). The
\textbf{improvement test lives in \texttt{validate}}, before any safety
check is evaluated: the \texttt{select\_and\_validate} loop
(\texttt{app/agents/workflow.py}) skips any candidate whose effective
switch delay is not strictly less than the current path's, records each
rejection with \texttt{failed\_checks = [not\_an\_improvement]}, and never
sends such a candidate to \texttt{validate\_reroute}. Only candidates that
survive the pre-filter are put through the nine named checks of Section~V-C.

Consequences, now stated explicitly in the paper:

\begin{itemize}
  \item The nine safety checks never see an improvement test; they are pure
    named predicates. \texttt{not\_an\_improvement} is a record tag, not a
    tenth check.
  \item If every recorded rejection is a non-improvement, the run reports
    \texttt{no\_better\_path} (S4; 5-of-5 trials in the earlier evaluation).
  \item If the candidate set is empty on entry to \texttt{validate} (the
    topology offers no alternative), the run reports \texttt{action\_rejected}
    with ``No usable alternative path exists in the current topology,''
    \emph{not} \texttt{no\_better\_path}; Algorithm~1 Line 22's old
    ``if every rejection failed only the improvement test'' branch is gone.
  \item The empty-rejection-set corner is no longer vacuously readable as a
    refusal; the outcome depends on which of the two situations the run is
    in, and the paper says which.
\end{itemize}
\end{response}

\subsection{Packet-loss clause in the violation predicate (Eq.~4) vs.\ the
loss target predicate in verification (Eq.~7)}

\begin{reviewer}
Because $v(f^+,I)$ in Eq.~(4) already returns true whenever $p_{f^+} > 0$,
any post-action packet loss forces $\neg v(f^+,I)$ to be false.
\end{reviewer}

\begin{response}
Confirmed, and fixed at the correct layer. The loss branch did not just
mis-order the verdicts---it did not exist in the code at all:
\texttt{safety.gates.latency\_violation} (shared by both gates) evaluates
the latency distribution only. The paper's Eq.~(4) therefore carried a
clause with no implementation, and it made Eq.~(7)'s third row unreachable
on paper. Fixed by construction:

\begin{itemize}
  \item Eq.~(4) is now exactly the code's predicate: unknown / violated /
    not-violated over the latency sample distribution, with the fallbacks for
    $n<2$ unchanged.
  \item The loss target is defined as the boolean $L \equiv
    (\ell \text{ undeclared}) \vee (p_{f^+} \le \ell)$. $L$ is the only
    place packet loss enters the verdict, and it is evaluated only by
    Eq.~(7).
  \item $Q$ is defined as the conjunction over the protected-flow
    evaluations $\mathcal{Q}$ (vacuously true when none), and ``ok'' is
    defined per flow: measured latency within its own constraint, loss within
    its own limit (if constrained), and path unmodified.
  \item $a$ and $b$ are defined as the scalar mean round-trip times
    $\bar{s}(f^+)$ and $\bar{s}(f^-)$ that the lab reports.
  \item It is now stated explicitly that $\neg v(f^+,I)$ in rows 2 and 3 of
    Eq.~(7) reads only the latency branches of Eq.~(4), and that $L$ carries
    the loss constraint. Row 3 (latency satisfies the target, loss does not)
    is reachable, exactly as the code has it.
  \item §V-D now tabulates the 1-to-1 mapping between Eq.~(7)'s verdict
    labels and the recorded outcomes: unverified $\rightarrow$
    \texttt{unverified}; success $\rightarrow$ \texttt{intent\_satisfied};
    partial $\rightarrow$ \texttt{partial\_improvement}; failed
    $\rightarrow$ \texttt{recovery\_failed}. No evaluated run produced
    \texttt{unverified}; it is kept distinct rather than collapsed into
    \texttt{recovery\_failed}.
\end{itemize}
\end{response}

\subsection{Suspect-link selection ($C_{\mathrm{susp}}$) and ordering of
freshness vs.\ violation checks in Gate 1 (\texttt{evidence})}

\begin{reviewer}
The gate names a suspect link \dots (Pipeline)\dots $C_{\mathrm{susp}}$ is a
set of links \dots In the execution trace of Scenario S2 \dots names both
\texttt{s2-s4} and \texttt{s1-s2}\dots $A(p)$ disqualifies any candidate
that shares that healthy upstream link (such as the four-hop path).
\end{reviewer}

\begin{response}
Two changes, one editorial and one in the paper's prose:

First, the rule is now stated exactly, and Section~VI-B spells it out:

\begin{enumerate}
  \item utilisation rule (fires first): every link with
    $u_e \ge 80\%$ of capacity is named, most-utilised first;
  \item path-delay fallback (no congested link): \emph{all} links on the
    flow's current path are named, sorted by effective delay descending,
    ties broken by link id. This is a set, and it is why S2's recorded
    suspects are $\{s2\text{-}s4 (48~\mathrm{ms}\ \mathrm{eff.}),\ s1\text{-}s2
    (4~\mathrm{ms})\}$, in that order---two links, not a single mis-flagged
    one. Equally, it is why a candidate such as $s1$-$s2$-$s3$-$s4$, which
    avoids the bottleneck entirely but shares $s1$-$s2$ with the current
    path, is disqualified by Eq.~(5). We do not present this as a feature;
    it is the cost of a conservative path-delay diagnosis, and S8's success
    depends on the identical mechanism.
\end{enumerate}

Second, the per-link configured delays are now given explicitly in
Section~VII-A ($d_{s1\text{-}s2}=4$, $d_{s2\text{-}s4}=4$,
$d_{s1\text{-}s3}=5$, $d_{s2\text{-}s3}=4$, $d_{s3\text{-}s4}=7$~ms, access
edges $0.1$~ms each), and the arithmetic is traceable end to end: healthy
RTT $\approx 2(8 + 2{\times}0.1) + 0.9 = 17.3$~ms, matching the
measurement.

On the freshness ordering: the paper's Algorithm~1 and Section~VI-B now say
explicitly that the routing consumes the feasibility precondition first. The
routing function (\texttt{app/agents/graph.py::\_evidence\_route}) was also
incorrect in the same way the old pseudocode was: it returned
\texttt{satisfied} when the samples satisfied $\tau$ even if the snapshot
was stale. That was a real defect (fresh evidence was never a precondition
of the confident no-op). Fixed:

\begin{itemize}
  \item violated $\wedge$ sufficient (fresh, sane) evidence $\rightarrow$
    diagnose;
  \item violated $\wedge$ insufficient evidence $\rightarrow$ abstain;
  \item not violated $\wedge$ sufficient evidence $\rightarrow$ satisfied;
  \item not violated $\wedge$ insufficient/stale/insane evidence
    $\rightarrow$ abstain --- never satisfied.
\end{itemize}

The fix is covered by a new regression test,
\texttt{test\_a\_stale\_snapshot\_is\_not\_reported\_as\_satisfied} (a stale
verdict routes to abstain; a sufficient one to satisfied; a violated one to
diagnose), and the routing docstring and Algorithm~1 lines are updated to
match.
\end{response}

\subsection{Figure 2 vs.\ thirteen compiled nodes, early exits and the
\texttt{awaiting\_approval} condition}

\begin{reviewer}
Figure 2 depicts 14 boxes \dots Algorithm 1 begins with guardrail \dots
omitted from Figure 2. \dots how early-exit runs execute the recording
step at Line 34. \dots The condition under which an action that passes all
nine checks still requires human approval \dots is not defined.
\end{reviewer}

\begin{response}
\begin{itemize}
  \item The mapping from the 13 compiled nodes to the boxes of
    Fig. 2 (the workflow figure) is now stated explicitly in Section~VI-A and in
    the figure caption: \texttt{intent}, \texttt{observe}, \texttt{evidence},
    \texttt{diagnose}, \texttt{validate}, \texttt{execute},
    \texttt{observe\_after}, \texttt{verify} plus five terminal nodes
    (\texttt{abstain}, \texttt{satisfied}, \texttt{stop}, \texttt{hold},
    \texttt{finish}). The guardrail is not a separate node; it runs inside
    \texttt{intent} before parsing and is what produces
    \texttt{guardrail\_blocked}; the ambiguous-intent clarification shares
    \texttt{stop}. The caption of Fig. 2 (the workflow figure) now says the same.
  \item \texttt{finish} runs only after \texttt{verify}; every early exit
    terminates through its own terminal node wired to END and does not pass
    through \texttt{finish}. The API persists the run record (outcome, gate
    verdicts, measurements) from the terminal state regardless of exit.
    Algorithm~1 was rewritten to state this instead of ending with an
    unconditioned ``finish'' line.
  \item The \texttt{lab\_offline} preservation is stated under both
    Section~VI-A and Algorithm~1: \texttt{observe} failing routes to the
    \texttt{abstain} terminal, and the \texttt{lab\_offline} flag is not
    overwritten.
  \item The risk-tier rule is now stated as it is implemented: HIGH when a
    protected-flow or operation-allowlist check failed; MEDIUM when a
    capacity, path-well-formedness or flow-scope check failed; otherwise
    LOW. A \emph{failed} check routes to rejection, never to execution. The
    hold path (\texttt{awaiting\_approval}) exists for an \emph{approved}
    action whose tier is in the unattended-hold set (default MEDIUM, HIGH);
    but in the evaluated configuration an approved action passes every
    check, so its tier is LOW and the hold path is never taken. The paper
    says this plainly instead of implying a hold happens on a failing check
    --- it never does.
\end{itemize}
\end{response}

\subsection{Exclusion of $C_{\mathrm{susp}}$ and effective delay from Gate
2's capacity check (\texttt{capacity\_respected})}

\begin{reviewer}
Gate 2's nine checks would admit [a path across an unflagged delay-degraded
link]. \dots under an unflagged pure-delay fault \dots Gate 2 does not
independently reject a candidate path traversing the degraded bottleneck.
\end{reviewer}

\begin{response}
Conceded, and now stated three times, at the point of use:

\begin{itemize}
  \item Section~V-C (Check 6) adds a paragraph: Check~6 tests
    $e\notin C_{\mathrm{deg}}$ and $u_e < 80\%$ only; it does not read
    $C_{\mathrm{susp}}$ and applies no effective-delay threshold. In S8
    ($C_{\mathrm{deg}}=\varnothing$, $u_e \approx 0$) Check~6 passes
    vacuously for every candidate, including straight across the fault; the
    repair succeeds only because the planner's Eq.~(5) exclusion removed the
    bottleneck candidate before Gate~2 was consulted.
  \item Section~VI-C (Gate 2) carries the same admission, in the same
    words.
  \item Section~VIII (scenario S8) reports the verdict with this caveat
    attached to the number, and Section~X (Limitations, ``Gate 2's capacity
    check is weaker than the planner's avoidance rule'') lists the decoupling
    as an open boundary; roadmap item~4 (Section~XI) is now precisely that
    decoupling.
\end{itemize}

We did not silently strengthen Check~6 with an effective-delay test in this
round: the change would alter the evaluated artifact's behaviour without a
rerun-verified migration, and per the review's own citation expectations the
honest move was to scope the claim, which the revised paper does. See
Section~4.1 for the S7/S8 utilisation numbers that make the vacuous case
empirically visible.
\end{response}

\subsection{Minor corrections and typos (system model and gates)}

\begin{response}
\begin{itemize}
  \item The uncompiled \texttt{\textbackslash emphdiagnosis} artifact in
    Section~IV-B (it rendered as ``emphdiagnosis'') is fixed---the intended
    \texttt{\textbackslash emph{diagnosis}} now compiles as
    \emph{diagnosis}.
  \item Utilisation is denoted $u_e$ uniformly (Section~IV-B, Eq.~(6),
    Section~V-C Check 6, Section~VIII).
  \item Check 1 now refers to $\mathrm{op}(a)$ on an action explicitly
    written $a=(\mathrm{op}, s, t, p)$, with \texttt{op} the constant
    \texttt{reroute\_flow} of the evaluated system.
  \item Path nodes are indexed $p = (v_0, \dots, v_k)$, $v_0 = s$,
    $v_k = t$ consistently (Section~IV-A defines it; Section~V-C Check 3
    uses it).
  \item The unused $V(p)$ / $V_{\mathrm{prot}}(f)$ preamble in Section~V-C is
    removed; the prose now explains that an earlier node-set comparison was
    replaced by the edge-set check $E_p \cap E_{\mathrm{prot}} =
    \varnothing$.
  \item ``The second row is a distinct outcome \dots'' is corrected to
    ``The third row \dots'', matching Eq.~(7)'s actual layout.
\end{itemize}
\end{response}

\section{Experimental setup, live-lab evaluation and ablations}

\subsection{S7: the fault was confounded and the utilisation was
sub-threshold}

\begin{reviewer}
The background UDP stream on the $100$~Mbit/s bottleneck link reaches
60~Mbit/s and registers as $\approx 56\%$ measured utilization \dots both
Gate 1's suspect-link identification and Gate 2's capacity check remain
unexercised by byte counters \dots violation detection, bottleneck
localization and path disqualification in S7 are driven by the exact same
\texttt{netem} delay and $C_{\mathrm{deg}}$ flag branches already tested in
S2 \dots
\end{reviewer}

\begin{response}
S7 now isolates the $\ge 80\%$ utilization branches for real, in the live
lab, with no lab-side oracle:

\begin{itemize}
  \item \texttt{s2-s4} is reconfigured, per trial, to a configured
    capacity of $60$~Mbit/s (the lab's \texttt{set\_link\_capacity}, which
    feeds the telemetry collector's denominator and the topology view; the
    Mininet link itself is unchanged). This models the emulated link rate
    a physical deployment would have in its monitoring configuration; the
    byte counters measured against it are real.
  \item A background UDP stream is offered at $70$~Mbit/s from
    $h2 \rightarrow h4$. Measured utilisation on the bottleneck is
    $80$--$94\%$ of the reconfigured capacity (mean $91\%$, $\approx 56$~Mbit/s
    of forwarding, host-CPU-bound---the userspace datapath's ceiling, stated
    plainly), and it crosses Gate 1's suspect threshold and Gate 2's
    Check~6 capacity threshold on the datapath's own byte counters in every
    sample.
  \item No lab-side degraded flag is set in S7
    (\texttt{report\_degraded=False}), so $C_{\mathrm{deg}}=\varnothing$
    and the localization must come from telemetry: the utilisation rule
    fires first and names $s2$-$s4$ (Section~VI-B).
  \item The $40$~ms netem injection remains in S7 \emph{only} as the fault
    that breaches the video flow's 50~ms intent---without it the flow sits
    at $\approx 17$~ms and the scenario would be a no-op---and its
    removal from the utilisation and flag paths is the point: utilization
    branches are exercised by the background traffic itself.
  \item A new lab endpoint \texttt{/v1/scenarios/link\_capacity} was added
    (\texttt{run\_lab.py}), proxied by the backend
    (\texttt{lab\_bridge.py}, \texttt{routes/lab.py}) and invoked by the
    harness for S7 only; every other scenario runs at the default
    configuration because the harness restarts the lab per scenario.
\end{itemize}

Measured outcomes (5 independent trials, each from a freshly built lab):
the bottleneck registers $\approx 56$~Mbit/s at $80$--$94\%$ utilisation
(mean $91\%$);
the video flow breaches the intent ($97$--$117$~ms pre-action during the
run), and the agent reroutes to $s1$-$s3$-$s4$; the verification gate
reports \texttt{intent\_satisfied} in every trial (post-action mean
$25.1$--$30.4$~ms; trial 2 was inflated by host scheduling). A tight
post-action distribution is reported in the paper as the trials'
min--max range, not as a single draw.

What S7 does not show, stated plainly: queueing behaviour, packet loss or
saturation of hardware line rate. The $80\%$ gate branches are now exercised
with live byte counters, but on a deliberately narrowed $60$~Mbit/s
configuration; Section~X (Limitations) says exactly this, and roadmap
item~3 pushes the rate past the narrowed capacity and the $100$~Mbit/s
bottleneck in hardware.
\end{response}

\subsection{Numerical discrepancies between prose and Tables II / VI for
peak latencies, and direction of disagreement in Ablation 2}

\begin{reviewer}
Section VIII-A prose states for S2 that the latency rose to 97.6 ms (peak
103.1 ms). However, Table II reports Peak = 138.8 ms \ldots and Section IX-C
prose states \dots ``or the reverse'', which the single row does not show.
\end{reviewer}

\begin{response}
All three were fixed, and the values come from the data pipeline rather
than from hand-editing:

\begin{itemize}
  \item The prose now quotes the same numbers the tables print: the peak
    is the largest single ping sample of the phase (17.5~ms for S1;
    138.8~ms for S2 in the last evaluated run, from a host-scheduling
    outlier); \texttt{numbers.tex} is generated by
    \texttt{make\_paper\_tables.py} from \texttt{data/}, so the prose and
    the tables cannot disagree.
  \item Table VI (the violation-rule ablation) now contains both directions, one row each:
    \texttt{no / yes} (outlier tolerance: one 138.8~ms spike pulls the
    10-sample mean above the target while fewer than $\gamma$ of the
    samples exceed it) and \texttt{yes / no} (oscillation: the majority of
    samples exceed the target while the mean sits just below it). The
    columns are explicitly defined in the caption: system / mean-only =
    each rule's verdict on the intent (yes = violated). The rows are chosen
    per direction as the closest calls (smallest $|\bar{s}-\tau|$) from a
    millisecond sweep of the target from 5 to 200~ms---the same boundary
    probe the \texttt{tab:ranking} ablation uses---not an estimate of the
    operating point. The $50$~ms intent of every other scenario remains the
    operating point, and no disagreement exists at it.
  \item Section~IX-C's prose now states both directions in the terms the
    data shows, and is specific about the 103.1~ms mean: it is the
    10-sample mean of the perturbed snapshot containing the
    138.8~ms host-scheduling outlier, while the other nine samples remain
    below the target---which is what makes the outlier tolerance the
    interesting case rather than the reverse.
\end{itemize}
\end{response}

\subsection{S8: where \texttt{d\_e\textsuperscript{eff}} comes from}

\begin{reviewer}
If the lab's snapshot API in S8 still directly populates the injected
40 ms netem delay inside the per-link telemetry tuple while merely omitting
\texttt{s2-s4} from the boolean set $C_{\mathrm{deg}}$, the snapshot is
still providing a link-level delay oracle rather than inferring link delay
from end-to-end measurements.
\end{reviewer}

\begin{response}
Correct, and it is now said in the paper rather than left for the reviewer
to puzzle out. Per-link $d^{\mathrm{eff}}_e$ in this lab is populated by the
lab itself: the netem injection mechanism (same as the S2/S8 fault) writes
the delay the link is actually carrying, the aggregator reports it, and S8
removes only the boolean \texttt{degraded} flag and the congested status.
It is not inferred from end-to-end probes---the lab cannot and does not
do that today. Section~VII-B states exactly what the snapshot supplies,
Section~VI-B describes the suspect rule the gate applies to it, and
Section~X scopes the claim: S8 demonstrates that Gate~1 localises from
the delay telemetry it is given, not that delays can be inferred from
end-to-end probes; roadmap item~4 now covers exactly that (including the
Gate~2 decoupling). A deployment that had no per-link delay telemetry
today—only end-to-end probes—has not been demonstrated to work, and this
paper does not claim it.
\end{response}

\subsection{Single-trial evaluations S3, S7, S8; unexercised Gate 2 rejection
paths; omission of control-loop execution time}

\begin{reviewer}
Repeat S3, S7, and S8 for 5 independent trials \dots include an ablation or
test case in Section IX where unsafe candidate paths are passed to Gate 2 to
empirically demonstrate \texttt{action\_rejected}; and briefly report the
end-to-end wall-clock execution duration of a recovery run in Section~VIII-A.
\end{reviewer}

\begin{response}
S3, S7 and S8 are now repeated as five independent trials each (S2 already
was), from freshly built labs, and their outcome distributions appear in
Table III (the outcomes table) and Table IV (the aggregates table): S3 reports
\texttt{partial\_improvement} 5-of-5 (post-action 25.4--26.1~ms); S7
reports \texttt{intent\_satisfied} 5-of-5; S8 reports
\texttt{intent\_satisfied} 5-of-5.

The wall-clock duration of the whole closed loop is now recorded per run
(\texttt{wall\_clock\_s}) and reported in Section~VIII-A: S2 mean
$\approx 5.8$~s, S3 $\approx 5.9$~s, S7 $\approx 6.3$~s and S8 $\approx 6.0$~s,
measured from the first observation to the verified outcome. Gate~1's five-second
freshness policy does not bound the loop as a whole; it bounds each
individual measurement, and the paper says so explicitly rather than
implying the loop fits inside five seconds.

On ``include an ablation or test case in Section~IX where unsafe candidate
paths are passed to Gate~2 to empirically demonstrate
\texttt{action\_rejected}'': we did not. The nine checks of Gate~2 are
unit-tested individually and as a conjunction, and the gate's rejection path
is exercised directly (e.g. \texttt{test\_validator\_rejects\_a\_reroute},
\texttt{test\_protected\_flow\_path\_is\_rejected},
\texttt{test\_wrong\_flow\_is\_rejected}), but no
scenario in the evaluated configuration proposes an inadmissible candidate,
because the planner's Eq.~(5) enumeration excludes them upstream of Gate~2.
\texttt{action\_rejected} is reachable and tested, not naturally exercised;
in the paper, Section~IX now says exactly that, and the outcome remains an
unexercised-in-data but tested branch, labelled as such (Section~X).
\end{response}

\subsection{Ambiguity in candidate scoring phrasing (Ablation 3)}

\begin{reviewer}
An always-act policy on a healthy network ``would have installed
s1-s3-s4---the best-scoring alternative, scoring better by Eq.~(6) since
nothing was wrong with the network.'' Under Eq.~(6), lower scores are
better, and on a healthy network at $0\%$ utilization, $s1$-$s3$-$s4$
actually scores \emph{worse} (higher) than the active path.
\end{reviewer}

\begin{response}
Correct and precisely the kind of inconsistency this round was for. Eq.~(6)
ranks by total delay first; $12$~ms $> 8$~ms, so $s1$-$s3$-$s4$ scores
worse than the active path. Section~IX-D now reads: the always-act variant
would have installed $s1$-$s3$-$s4$---the best-scoring \emph{remaining
alternative after excluding the current path}, selected because a variant
that bypasses the strict delay-improvement filter selects it (rather than
because it outranks $s1$-$s2$-$s4$). The Table (\texttt{alwaysact})'s
``always act'' row and its measured $+50\%$ regression are unchanged; only
the reasoning was corrected.
\end{response}

\subsection{Minor corrections (experimental setup)}

\begin{response}
\begin{itemize}
  \item The phases table's caption and the latency-phases and
    sample-distribution figures now state explicitly that S6 (lab stopped)
    yields no samples and is omitted from both the table and the axes, while
    it appears in Table III (the outcomes table) with its
    \texttt{lab\_offline} outcome.
  \item ``a $8$~ms path'' is now ``an $8$~ms path''.
  \item ``An undegraded healthy network will ought to produce" is now
    rewritten as part of the S7 paragraph (Section~VIII-A) and no longer
    contains the double modal.
  \item The scenario tick labels on the latency-phases and
    sample-distribution figures are shortened (\texttt{S2 $\times$5},
    \texttt{S3 $\times$5}, \texttt{S7 $\times$5}, \texttt{S8 $\times$5}) and
    rotated $30^{\circ}$; the phases' x-axis labels no longer collide.
\end{itemize}
\end{response}

\section{Limitations, audit takeaways and the future roadmap}

\subsection{Internal contradictions between the roadmap and earlier sections}

\begin{reviewer}
Section XI appears to have been carried over from an earlier revision of
the manuscript \dots Roadmap Item 2 (``Enforce the rules in the data
plane'') \ldots However, Section VII-A states the matcher fix \dots Roadmap
Item 4 (``Decouple detection from injection'') \ldots S8 has already been
implemented and evaluated \ldots Roadmap Item 5 (``Fix the reset
semantics'') \ldots Section X states the reset now restores the baseline
forwarding path.
\end{reviewer}

\begin{response}
All three were rewritten; the Abstract no longer says the reset does not
restore the datapath's learned path either.

\begin{itemize}
  \item Item 2, now ``Port the evaluation to a real switch path'': the
    user-space matcher bug is fixed and covered by regression tests; the
    remaining work is porting the evaluation to Open vSwitch (userspace
    mode) and then to hardware, eliminating Python interpreter scheduling
    jitter and the $\sim$56~Mbit/s forwarding ceiling, and seeing
    $c_{\mathrm{dp}}$ fall toward a merchant-silicon value.
  \item Item 4, now ``Decouple Gate 2's capacity check from $C_{\mathrm{deg}}$'':
    S8 (this paper) removed the lab-side oracle flag and Gate~1 localised
    the bottleneck from effective-delay telemetry alone; the remaining
    dependency is Gate~2's Check~6, which tests $e\notin C_{\mathrm{deg}}$
    and $u_e$ only and passes vacuously in the unflagged case (Section~V-C,
    Section~VI-C, Section~X). The decoupling---an effective-delay-based
    capacity predicate in Gate~2, or inferring link delays from end-to-end
    probes---is the remaining work, and it is now scoped that way.
  \item Item 5, now ``Compound the fix, then keep the restart'': the claim
    that the harness restarts because the reset is broken is obsolete;
    the reset restores the baseline forwarding path (measured: 25.4~ms
    after the reroute, 17.3~ms after reset). The restart is retained only
    because a bit-identical baseline per trial costs a few seconds; retiring
    it is an operational choice, not an open defect. The Abstract's
    sentence to the contrary is gone.
\end{itemize}
The Abstract, Section~X and Section~XI are now mutually consistent. We
checked every roadmap item against the completed sections line by line.
\end{response}

\subsection{Multi-flow interactions and standardized benchmarking in the
roadmap}

\begin{reviewer}
When scaling beyond a single controlled flow, the authors might also
address how the deterministic planner and safety gate will account for
multi-flow cascading effects \ldots For Roadmap Items 1 and 6, the
execution plan could be further grounded by referencing established
benchmarks in this domain, such as NetConfEval or NIKA.
\end{reviewer}

\begin{response}
Both were incorporated. Roadmap item~6 now proposes not only porting the
scenarios to a larger topology but also accounting for multi-flow cascading
effects: when rerouting $f_1$ shifts offered load onto an alternative
path, a non-protected flow $f_2$ may be pushed over the $80\%$ utilization
threshold or past its own latency objective---so the planner and Check~6
must account for the counterfactual load, and Gate~2's protection semantics
(Check~7) must be extended from a single protected flow to a set of them
competing for the same edges. Roadmap item~1 now suggests evaluating the
active model against \textbf{NetConfEval} (Changjie Wang et al., 2024) for
natural-language intent specification and conflict detection, and
\textbf{NIKA} (Zhihao Wang et al., 2025) for benchmarking troubleshooting
agents across diverse Mininet/SDN topologies and fault categories, so the
boundary is measured against published, reproducible tasks rather than only
our own four-switch lab.
\end{response}

\subsection{Minor corrections (roadmap and bibliography)}

\begin{response}
\begin{itemize}
  \item Reference [6] (RFC 9315) no longer carries a duplicated year; it
    reads ``RFC 9315, October 2022'' once.
  \item Reference [12] (Mondal et al.) is now cited as the HotNets '23
    peer-reviewed paper: \emph{Proceedings of the 22nd ACM Workshop on Hot
    Topics in Networks (HotNets '23)}, pp.\ 189--195, 2023. (The DOI was
    deliberately omitted rather than guessed.)
\end{itemize}
\end{response}

\section{Citation and reference verification}

\begin{reviewer}
Paper References Checked: 13. Verified: 13. Flagged (Likely Hallucinated): 0.
Unresolved: 0.
\end{reviewer}

\begin{response}
The verification found no hallucinated references, and the two flagged
entries (both of which the review had singled out in its detailed reading)
are corrected as in Section~5.3. Five new references were added, each
checked against the publisher's own record: Sanvito et al.\ (NetSoft 2018,
DOI~10.1109/NETSOFT.2018.8460064, pp.~272--276), Jacobs et al.\ (USENIX
ATC '21, pp.~625--639), Wang et al.\ (NetConfEval, Proc.\ ACM Netw., 2024,
DOI~10.1145/3656296), Wang et al.\ (Confucius, SIGCOMM '25,
DOI~10.1145/3718958.3750537), Protogeros et al.\ (arXiv:2604.22513, 2026),
and Wang et al.\ (NIKA, arXiv:2512.16381, 2025).
\end{response}

\section{Verification: what was rebuilt, re-measured and re-passed}

\begin{itemize}
  \item \textbf{Reproduction.} \texttt{bash final\_submission/scripts/reproduce.sh}
    was re-run end-to-end with the fixes in this round; it rebuilt the lab,
    ran all eight scenarios (with S2, S3, S7, S8 as five independent trials
    each), re-ran the suites, and regenerated every figure and the paper's
    tables and numbers from the recorded data.
  \item \textbf{Measured results (this run).} Healthy 17.6~ms; degraded
    97.0--97.5~ms; recovery 25.4--25.6~ms across five trials (5/5
    \texttt{intent\_satisfied}); partial improvement 25.4--26.1~ms (5/5
    \texttt{partial\_improvement}); no-oracle 25.4--25.5~ms (5/5
    \texttt{intent\_satisfied}); offered load 25.1--30.4~ms (5/5
    \texttt{intent\_satisfied}), bottleneck at 60~Mbit/s with the
    utilisation branch exercised on live byte counters ($80$--$94\%$, mean
    $91\%$ of the reconfigured capacity, no degraded flag); always-act
    comparison: $17.6 \rightarrow 25.8$~ms ($+46\%$).
  \item \textbf{Wall-clock.} The closed loop completes in a few seconds per
    run (S2 $\approx 5.8$~s, S3 $\approx 5.9$~s, S7 $\approx 6.3$~s,
    S8 $\approx 6.0$~s), i.e.\ each individual measurement is captured and
    consumed within the five-second freshness policy's scope---the policy
    governs individual measurements, not the loop.
  \item \textbf{Test suites.} 134 backend tests pass (including the new
    stale-snapshot routing regression test), 4
    network-lab datapath tests pass, 5 live-lab tests pass.
  \item \textbf{Paper.} The paper rebuilt with \texttt{tectonic}; its
    tables and every macro quoted in its prose are generated by
    \texttt{make\_paper\_tables.py} from the recorded data (nothing is
    hand-typed), and every reference was verified against the publisher's
    record.
\end{itemize}

\end{document}
